\section{Conclusion}\label{sec:conclusion}

We have presented a numerically stable framework for decomposing pretrained ReLU networks into split-pairs of monotone and convex split-streams, enabling improved explainability without retraining or altering model predictions.
While there is no single saliency method excellent in \emph{all} metrics at the same time, our experiments establish splitting as a tool to boost various of them on trained networks. 
By directly training DC-models, we initiate their application as a further approach to increase explainability of a network. 
Especially the separation between factual and counterfactual attributes observed in the MNIST experiments is remarkable and may be understood further theoretically and experimentally.

While we provided some details on the splitting of feedforward networks, with similar ideas, this can be extended to further architectures. 
We provide a novel mathematical framework to decompose multiplicative networks as the transformer \citep{Vaswani2017} into the difference of monotone and convex parts in \cref{app:splitting_multiplicative_maxout_networks}.
It is future work to analyze the properties of the proposed decomposition approach for such multiplicative networks and to adapt SplitCAM and SplitLRP to the transformer architecture.

The split network need not duplicate parameters: instead of materializing a separate model, one can traverse the original \texttt{nn.Module} graph and realize the split using masked views of the same weights. It is future work to implement the network decomposition framework more memory efficient without duplicating model weights.

Our experiments show that attribution quality heavily depends on layer selection, with different layers excelling at different metrics. 
This motivates to develop principled methods for automatic layer selection or combining attributions across multiple layers for more robust explanations.